% =============================================================================
%  CHAPTER 8 -- NONLINEAR MODELS, OVERFITTING, AND THE RESET TEST
% =============================================================================
\section{Nonlinear models, overfitting, and the RESET test}
\label{ch:nonlinear}

Phase~\ref{ch:phase9} asks a specific question: does a flexible nonlinear forecaster beat
the regularized linear map on the same causal features, out of sample, net of cost? This
chapter builds the theory needed to answer it and to predict the answer in advance.

\subsection{The estimation target, and how to score a forecast}
\label{sec:nonlinear:target}

Everything is a forecast of a conditional mean. Write the \emph{regression function}
\begin{equation}
  m(f) \;=\; \E\bigl[R \mid F=f\bigr],
  \label{eq:regressionfunction}
\end{equation}
where $F_t$ is a causal feature vector (in our case 11 features built from bars
$\le t$: the current $z$, the spread level and its changes, 5- and 20-bar rolling
volatility of the spread, target and basket returns and their lags) and $R_t$ the target
(the realized $h$-bar market-neutral return of fading the spread,
equation~\eqref{eq:nonlin_target}). Any estimator $\hat m$ is judged by its
out-of-sample mean squared error, which decomposes exactly:
\begin{equation}
  \E\bigl[(R-\hat m(F))^2\bigr]
  \;=\; \underbrace{\E\bigl[(m(F)-\hat m(F))^2\bigr]}_{\text{estimation error}}
  \;+\; \underbrace{\E\bigl[(R-m(F))^2\bigr]}_{\text{irreducible noise } \sigma^2} .
  \label{eq:mse:decomp}
\end{equation}
(The cross term vanishes because $\E[R-m(F)\mid F]=0$ by definition of the conditional
mean.) Further decomposing the estimation error at a fixed $f$ over repeated samples,
\[
  \E_{\mathcal D}\bigl[(m(f)-\hat m(f))^2\bigr]
  = \bigl(m(f)-\E_{\mathcal D}\hat m(f)\bigr)^2 + \Var_{\mathcal D}\bigl(\hat m(f)\bigr)
  = \mathrm{bias}^2 + \mathrm{variance},
\]
which is \eqref{eq:biasvariance} applied to a function. The only lever is the
bias/variance split; $\sigma^2$ is a property of the market.

\paragraph{Two scores, and why we report both.}
\begin{itemize}
  \item \textbf{Out-of-sample $R^2$}: $1-\mathrm{MSE}_{\text{oos}}/\Var(R_{\text{oos}})$.
    It is scale-free but brutally sensitive to the noise term: with a daily reversion
    signal whose explained variance is a few percent, an $R^2$ of $0.05$ is a strong
    result and a negative $R^2$ merely means ``worse than predicting the mean''.
  \item \textbf{Information coefficient (IC)}: $\Corr(\hat m(F_t), R_t)$ out of sample.
    It ignores the forecast's scale and measures only directional/rank content, which is
    what a threshold policy consumes. Since $R^2 = \mathrm{IC}^2$ for a correctly scaled
    linear forecast, an IC of $0.34$ corresponds to $R^2\approx0.12$ --- and our measured
    linear IC of $0.341$ with $R^2=0.088$ (Table~\ref{tab:nonlinear_forecast}) shows the
    forecast is slightly \emph{mis-scaled} but well-directed. We also report a
    rank/persistence reference (IC of $z_t$ itself) so that ``beyond just trading the
    current mispricing'' is measurable.
\end{itemize}

\subsection{Tree ensembles: how they fit, and why they overfit}
\label{sec:nonlinear:trees}

\paragraph{One tree.} A regression tree partitions the feature space into axis-aligned
boxes by recursively choosing a split (feature $j$, threshold $s$) that minimises the
within-child sum of squares,
\begin{equation}
  \min_{j,s}\ \Bigl[\ \min_{c_L}\sum_{i\in L}(y_i-c_L)^2 + \min_{c_R}\sum_{i\in R}(y_i-c_R)^2\ \Bigr],
  \label{eq:tree:split}
\end{equation}
whose inner solutions are the child means, so the criterion reduces to maximising
$\frac{(\sum_{i\in L}y_i)^2}{\lvert L\rvert}+\frac{(\sum_{i\in R}y_i)^2}{\lvert R\rvert}$.
The prediction is piecewise \emph{constant}. Consequences: trees are invariant to
monotone transformations of features, handle interactions automatically, and
\emph{extrapolate terribly} --- outside the training range they return a boundary
constant, which for a spread signal means a large confident error precisely when the
spread does something unprecedented.

\paragraph{Boosting as functional gradient descent.} Gradient boosting fits an additive
model $F_M(x)=\sum_{m=1}^M \nu\,h_m(x)$ by, at each step, computing the negative gradient
of the loss at the current fit,
$\tilde y_i = -\bigl[\partial L(y_i, \hat y_i)/\partial \hat y_i\bigr]_{\hat y_i=F_{m-1}(x_i)}$
(which for squared loss is the residual $y_i-F_{m-1}(x_i)$), fitting a base tree to those
pseudo-targets, and adding it with learning rate $\nu$. With $\nu$ small the ensemble is
an explicit descent path in function space; with $\nu=1$ it interpolates the training
data quickly and variance explodes. \code{HistGradientBoostingRegressor} additionally
bins features into 255 buckets (histogram-based splits), which is both a speed trick and
an implicit regulariser.

\paragraph{Why it must lose here.} Boosting's bias is essentially zero given enough
rounds; all its error is variance, and variance scales roughly as
$O(\text{complexity}/T)$. With $T\approx2{,}900$ usable training bars, 11 features, and a
signal-to-noise ratio of a few percent, the variance term dominates. The measured
learning curve (Table~\ref{tab:nonlinear_learning},
Figure~\ref{fig:nonlinear_learning}) is the direct evidence: out-of-sample $R^2$ is
$-0.45$ at 400 training bars, $-0.13$ at 2{,}000, $-0.03$ at 2{,}400, and only at
$2{,}800$ bars does it turn positive ($+0.43$ on that particular validation window). The
model is still climbing out of its own variance hole at the very end of the available
data. There is no version of this dataset where it wins.

\subsection{The shallow network}
\label{sec:nonlinear:mlp}

A single hidden layer with $H$ units computes
\begin{equation}
  \hat m(f) \;=\; b_0 + \sum_{h=1}^{H} c_h\,\phi\bigl(w_h^{\top}f + a_h\bigr),
  \label{eq:mlp}
\end{equation}
with $\phi$ a nonlinearity (ReLU, tanh). The universal approximation theorem says such a
network can approximate any continuous function on a compact set arbitrarily well as
$H\to\infty$ --- a statement about \emph{bias} that says nothing about variance and is
therefore almost useless for a $T=3{,}000$ problem. Training is non-convex
(gradient descent from random initialisation, so results vary with the seed), the
effective capacity is controlled by early stopping and weight decay rather than by $H$,
and the whole apparatus is sensitive to feature scaling --- which is why the pipeline
standardises features using training-window statistics only (a leak that is easy to make
and is explicitly avoided here). Measured: the MLP's IC ($0.267$ on ADBE) sits between
the linear ($0.341$) and the GBM ($0.255$), and its $R^2$ is negative ($-0.049$).

\subsection{Ramsey's RESET test: is the linear form wrong?}
\label{sec:nonlinear:reset}

RESET (Regression Equation Specification Error Test, \cite{ramsey1969}) is a
general specification test that requires no alternative model. Given the fitted linear
model $y=X\beta+\eps$ with fitted values $\hat y=X\hat\beta$:

\begin{enumerate}
  \item Form the \emph{augmented} regression
    $y = X\beta + \gamma_2\hat y^2 + \gamma_3\hat y^3 + \dots + \gamma_k \hat y^k + u$.
  \item Test $H_0:\gamma_2=\dots=\gamma_k=0$ with the standard $F$-statistic
    \eqref{eq:ftest}, with $k-1$ numerator and $T-n-k+1$ denominator degrees of freedom.
\end{enumerate}

\paragraph{Why powers of $\hat y$?} If the true relation is nonlinear in $X\beta$, then
$y\approx g(X\beta)$ and a Taylor expansion of $g$ around $\E[\hat y]$ is a polynomial in
$\hat y$; adding those powers should improve the fit. The test is therefore a
low-order polynomial check on the functional form, and it is powerful against the most
common alternatives (quadratic curvature, saturation). Its weakness: it is not powerful
against omitted \emph{variables} that are orthogonal to $\hat y$, and using $\hat y$
rather than $X\beta$ makes the augmented regressors correlated with the original ones,
which mildly distorts small-sample size.

\paragraph{A caution about running it out of sample.} RESET is a test about the
\emph{fitting} equation, so it is run on the training sample only
(Phase~\ref{ch:phase9}, $n_{\text{train}}=2{,}881$); running it on held-out data would
mix two different questions (misspecification versus generalisation).

\paragraph{Measured result.} Table~\ref{tab:nonlinear_reset}: ADBE $p=0.376$ (does
\emph{not} reject linearity), AVGO $p<10^{-3}$, GS $p=0.0023$, CAT $p<10^{-3}$ (all
reject). The reading is subtle and worth stating carefully: on three of four baskets the
linear forecasting equation \emph{is} statistically misspecified, yet the flexible models
fail to convert that misspecification into out-of-sample $R^2$ (AVGO GBM: $+0.119$ versus
linear $+0.173$; GS GBM: $-0.052$ versus linear $+0.078$; CAT GBM: $-0.058$ versus linear
$+0.041$). Detecting nonlinearity and \emph{exploiting} it are different problems, and
the second requires far more data than the first. With $T\approx3{,}000$ and a
few-percent signal, exploitation fails.

\subsection{From IC to Sharpe: the fundamental law, and why it over-promises}
\label{sec:nonlinear:fundamentallaw}

Grinold's fundamental law of active management states
\begin{equation}
  \mathrm{IR} \;=\; \mathrm{IC}\times\sqrt{\text{breadth}},
  \label{eq:fundamentallaw}
\end{equation}
where IR is the (annualised) information ratio --- for a market-neutral book, essentially
the Sharpe --- and breadth is the number of \emph{independent} forecasts per year. It
follows from \eqref{eq:sharpetstat}: summing $B$ independent bets each with mean
$\mu_b=\mathrm{IC}\cdot\sigma_b$ and standard deviation $\sigma_b$ gives a portfolio mean
$B\mu_b$ and volatility $\sqrt B\sigma_b$, hence ratio $\mathrm{IC}\sqrt B$.

Applying it naively to our numbers produces an absurdity, and diagnosing the absurdity is
instructive. With $\mathrm{IC}=0.34$ (ADBE linear) and $B=252$ daily forecasts,
\eqref{eq:fundamentallaw} gives $\mathrm{IR}=0.34\times\sqrt{252}=5.4$. The measured net
Sharpe is $0.05$. Five things break the law's assumptions, in decreasing order of
importance here:
\begin{enumerate}
  \item \textbf{The forecasts are not independent.} The target is the realized
    \emph{$h=5$-bar} (and, for the IC column of Table~\ref{tab:model_compare}, 21-bar)
    forward return, so consecutive forecasts overlap and share the same future. The
    effective number of independent bets per year is $\approx252/h$, i.e.\ $50$ for
    $h=5$ and $12$ for $h=21$, giving $\mathrm{IR}\approx0.34\sqrt{50}=2.4$ or
    $0.34\sqrt{12}=1.2$. Still too high.
  \item \textbf{Breadth is not forecasts, it is \emph{positions taken}.} The threshold
    policy trades only when $\lvert z_t\rvert\ge2$ --- about $198$ round trips in
    4{,}190 bars, i.e.\ $\approx12$ trades per year, with the book flat (or unchanged)
    most of the time. Using $B=12$: $\mathrm{IR}\approx0.34\sqrt{12}=1.2$.
  \item \textbf{IC is measured gross of cost and against a partially hedge-dependent
    target.} The realized return $R_t^h$ is itself computed from a \emph{rolling} hedge,
    so a forecast that predicts the hedge's own estimation error gets credit that a real
    book cannot collect. At 10\,bps one-way and $\approx12$ round trips/year, cost drag is
    $\approx12\times2\times0.001=2.4\%$/year on \$1 gross, versus a gross annual return of
    $3.0\%$ (Table~\ref{tab:backtest_base}) --- i.e.\ \textbf{costs consume 80\% of the
    gross edge}. This is the dominant term, and it is why net Sharpe is $0.05$ rather than
    $0.385$.
  \item \textbf{The transfer coefficient is below 1.} Constraints (threshold bands,
    one-bar execution delay, dollar neutrality, integer rebalancing) mean the position
    actually held is not $\propto$ the forecast. In Grinold's full form
    $\mathrm{IR}=\mathrm{IC}\times\sqrt{B}\times\mathrm{TC}$ with $\mathrm{TC}\le1$;
    here $\mathrm{TC}$ is small because the policy is a sign function of $z$, not a
    proportional bet.
  \item \textbf{IC is estimated with error.} With $\approx1{,}164$ overlapping
    observations, the standard error of an IC estimate is roughly
    $1/\sqrt{T_{\text{eff}}}\approx1/\sqrt{55}\approx0.13$ for $h=21$ --- comparable to
    the IC itself. The honest statement is ``IC is positive'', not ``IC is $0.34$''.
\end{enumerate}

\begin{remark}[Honesty: the arithmetic of a weak edge]
Putting (2) and (3) together gives the real picture, and it is worth internalising
because it explains every negative result in Part~II: a single-basket daily spread
strategy takes $\sim$12 independent bets a year with an IC of $0.2$--$0.3$, so its
\emph{gross} Sharpe ceiling is $\approx0.3\sqrt{12}\approx1.0$, and paying 10\,bps
round-trip on every bet removes most of it. The measured gross Sharpe of $0.385$ and net
of $0.05$ are entirely consistent with that arithmetic. To do better you need more
breadth (many more baskets --- which multiplies the hypothesis count $M$ and therefore
the correction in Chapter~\ref{ch:multipletesting}), lower cost (fewer, larger trades),
or a genuinely higher IC. Nothing in the data suggests the third is available at daily
frequency with public equity data.
\end{remark}

\subsection{What would make a nonlinear model worth it}
\label{sec:nonlinear:when}

The verdict ``linear wins'' is a statement about \emph{this} sample size, feature set, and
horizon --- not a universal claim. Nonlinear models earn their variance when: (i) the
signal-to-noise ratio is higher (intraday microstructure, order-flow features); (ii) the
sample is much larger (a cross-section of hundreds of baskets pooled, so $T\times N$
observations rather than $T$); (iii) the relationship has genuine structure the linear
form cannot express (option-like payoffs, hard regimes, interactions with liquidity);
(iv) the target is a classification with strong asymmetry (``will this spread break?''
rather than ``how much will it revert?''). Item (ii) is the cheapest to test and is the
first item of future work in Chapter~\ref{ch:conclusion}: pool all 195 names' baskets into one
supervised-learning problem with basket fixed effects, which raises the effective sample
by two orders of magnitude while keeping the per-basket economics.
